% S1 Appendix of the manuscript (manuscript/main.tex). Build: pdflatex S1_Appendix.tex (twice). Every number is read from the sealed reports
% (archive/study_a/RESULTS_REPORT.md via STUDY_A_SUMMARY.md, part_b/RESULTS_PART_B.md, part_b/RESULTS_E3C_MU_DIAL.md, part_c/RESULTS_PART_C.md,
% part_c/PART_C_LOG.md, part_d/RESULTS_PART_D.md, part_d/PART_D_LOG.md), the preregistrations and their Deviations sections, the seal manifests
% (FREEZE*.json) and, for the numbers labelled post hoc, posthoc/results (definitions in posthoc/README.md).
% Build: ../build.sh si  (or pdflatex twice in this folder).
\documentclass[10pt,letterpaper]{article}
\usepackage[margin=1in]{geometry}
\usepackage{amsmath,amssymb}
\usepackage[T1]{fontenc}
\usepackage{lmodern}
\usepackage{microtype}
\usepackage{array}
\usepackage{enumitem}
\usepackage[hidelinks]{hyperref}
\setlist{nosep,leftmargin=1.4em}
\newcommand{\dd}{\mathrm{d}}
\newcommand{\tol}{\varepsilon_{\mathrm{tol}}}
\newcolumntype{L}[1]{>{\raggedright\arraybackslash}p{#1}}
\setlength{\parindent}{0pt}
\setlength{\parskip}{0.5em}
\raggedright

\begin{document}

{\Large\textbf{S1 Appendix. Secondary methods and results}}

\smallskip
Supporting information for ``A single pre-run time step loses to adaptive integration in the Earth--Moon three-body problem; in Sundman time it
wins on most trajectories but not in total cost''. Section numbers below are those of this appendix; ``post hoc'' marks analyses of the sealed
ledgers defined after all studies had run (\texttt{posthoc/README.md} in the repository), and ``exploratory'' marks analyses outside every
preregistered test family. File paths refer to the deposited repository.

\tableofcontents

\section{Provenance of the preregistrations}
Each study was preregistered in a document whose text above a ``Deviations'' heading was hashed (SHA-256) together with every source file, the
frozen models and the datasets by a seal script. The confirmatory runner of each study refused to start a step if any hash differed and logged
each step once.

{\small
\begin{tabular}{|l|L{1.75in}|L{0.95in}|L{2.6in}|}
\hline
\textbf{Study} & \textbf{Document} & \textbf{Sealed (UTC)} & \textbf{Recorded state of the code} \\ \hline
A & \texttt{PREREGISTRATION.md} & 2026-10-23 23:00:16 & code hash only (\texttt{freeze\_snapshot.py}); seal committed as \texttt{912af65} \\ \hline
B & \texttt{part\_b/PREREGISTRATION\_B.md} & 2026-10-25 23:01:57 & git \texttt{a3ee296}; a first seal (\texttt{FREEZE\_B.v1.json}) was replaced before any confirmatory step (Section~\ref{sec:dev}) \\ \hline
C & \texttt{part\_c/PREREGISTRATION\_C.md} & 2026-10-27 01:37:44 & git \texttt{ea3b6df}, clean tree \\ \hline
D & \texttt{part\_d/PREREGISTRATION\_D.md} & 2026-10-28 21:37:15 & git \texttt{4586803} plus 23 uncommitted paths, among them the Study D code, recorded by hash only; 83 source files and 60 frozen files hashed \\ \hline
\end{tabular}}

\textbf{What can and cannot be verified.} The seals are local hash manifests in a git repository that was not mirrored to a public host during
the work. Commit dates are set by the author's machine, so the order of events (design, seal, confirmatory run) cannot be confirmed by a third
party from the repository alone; the hashes only show that the sealed files have not changed since the manifest was written. The Study D seal is
weaker than the others because the Study D code was not committed at the seal; its hashes sit in a file that was committed later. Nothing in
the repository indicates that a sealed file was changed, but these preregistrations should be read as internal, not as registered reports.
The deposit of the repository with a DOI (manuscript, Data availability) timestamps its state at deposit; it cannot backdate the seals.

\section{Dated deviations of every study}
\label{sec:dev}
Each item is quoted in substance from the Deviations section of the study's preregistration or from its log; dates are those recorded there.

\subsection{Study A}
\begin{itemize}
\item 2026-10-23, before the seal: the certification of the safety scales, the global fixed step and the adaptive tolerance was changed from a
fixed sequence starting at the smallest value to a Bonferroni rule over the grid (largest certified value). Runs that hit the step cap counted as
violations and made the smallest scale the worst one (leapfrog 2.7\% step-cap failures at scale 0.02 in a 300-row test), so the fixed sequence
could certify nothing. Also before the seal: a random rather than grouped train/held-out split; adaptive solves capped at 2,000,000 RHS
evaluations.
\item 2026-10-23, before the seal: the wording of H1 was corrected to the implemented cost basis (H1 charged only the FTLE feature; the model's
prediction cost is charged only in H1$'$).
\item 2026-10-23: \textbf{H1$'$a and H1$'$b were added after development data had been seen.} On 80 development trajectories the FTLE feature
alone cost more than a whole DOP853 solve on every one (median factor about 70), so H1 could not be supported as written; batched model inference
was measured at about 730 RHS-equivalents against a median DOP853 solve of about 254 evaluations, so H1$'$a was expected to fail. H1 itself was not
changed.
\item 2026-10-23: the regressor (gradient-boosted, quantile 0.1) was pinned for every command, because the earlier run had evaluated a different
model on the locked set than the one developed.
\end{itemize}

\subsection{Study B}
\begin{itemize}
\item 2026-10-25, before any confirmatory step and before any confirmation or lock outcome had been analysed: the first seal was replaced. The
runner wrote the lock-set outputs into the confirmation-set folder (so a later step could have read the wrong policy), and three presentation
functions were added to the figure script. The preregistration text was unchanged (same hash).
\item 2026-10-26, after every confirmatory step had run: the figure script was edited for legibility only.
\item 2026-10-26, after the confirmatory run and before any result existed: two exploratory studies (the lookahead study E3c and the mass-ratio
dial E7) were preregistered in the Deviations section, with expectations and an alternative stated in advance (Section~\ref{sec:studyb}).
\item 2026-10-26, after they ran: executed as written, except that intervals were not computed for the pinball loss and the 90th percentile of the
error, and two runs that failed on a missing import and a column name were rerun with no change to any computed quantity.
\end{itemize}

\subsection{Study C}
Decisions made during development, before the seal, that departed from the study plan (all recorded in \texttt{part\_c/PART\_C\_LOG.md}):
\begin{itemize}
\item 2026-10-26: \textbf{the overhead gate for model selection was raised from 10 to 20 RHS-equivalents}, after the timing measurements and after
reading the integration-only costs of the selection stage but before any total-cost ranking. No candidate met 10 (the floor was 10.7 for
state-only features, because the periapsis features alone cost 3.7--3.9). The eligible set was fixed by the timing table alone.
\item 2026-10-26: \textbf{the Sundman-time label missed its acceptance criterion}: 9 of 200 pilot trajectories (4.5\%; 95\% interval 2.1--8.4\%)
were non-monotone in the fictitious step on a dense sweep, against the plan's $<2\%$, and 2.0\% had a failing step inside the window the label
claims to certify. The study proceeded because the plan's go rule did not include this criterion (Section~\ref{sec:labels}).
\item 2026-10-26: the leapfrog arm was dropped from the physical-time policy after a development comparison (Section~\ref{sec:leapfrog}).
\item 2026-10-26: the quantile level of the Sundman-time model was chosen on the selection set (0.2; the plan's 0.1 had not been tuned).
\item 2026-10-26: the best abstention threshold of the physical-time policy ($\log_{10}\hat\tau < -1.0$) was \textbf{at the edge of its preregistered
grid}; a better threshold may exist and was not searched, and the combination of the step-over guard with abstention was not run.
\item The bisection shortcut of the fixed-sequence certificate was used for selection only; every confirmatory certificate used the exact rule.
\end{itemize}
After the seal:
\begin{itemize}
\item 2026-10-27, after every confirmatory step had run: one exploratory addition (an arm of E3c with the periapsis features at no lookahead cost),
written before its code existed, with its expectation.
\item 2026-10-27: a figure title was edited for legibility.
\end{itemize}

\subsection{Study D}
Before the seal (\texttt{part\_d/PART\_D\_LOG.md}, sections 4 and 5):
\begin{itemize}
\item The kernels of the compiled wall-clock comparison (D2) were \textbf{corrected twice in dry runs on the Study D selection set} before the seal
(Section~\ref{sec:d2}). Neither discarded kernel produced a confirmatory number.
\item The D5 analysis gained its cluster bootstrap and the D10 and IAS15 commands were written; the expectations of the preregistration were filled
in from the dry runs.
\item The seal recorded the commit plus the hashes of \textbf{23 uncommitted paths}, among them the Study D code (Section~1).
\end{itemize}
After the seal:
\begin{itemize}
\item 2026-10-28, before the D2 step: the 1-minute load average did not fall below the guard of 1.5 (it sat at 1.49--1.50 with the processor 94.5\%
idle), so the step was launched with the guard \textbf{relaxed to 1.6}. The load recorded at the start of the step was 1.499, below the original
guard, so the relaxation turned out not to be needed; it is recorded because the step was launched with it.
\item 2026-10-29, after every confirmatory step had run: the per-stratum certified adaptive comparator (D4b) was written \emph{after} the D4 result was
known and because of it; it is post hoc. A figure legend was moved.
\item The independent reference method of the exploratory check D12 was changed from Radau (which did not finish within 10 minutes at $10^{-13}$) to
LSODA before any use.
\end{itemize}

\subsection{Post hoc analyses}
The post hoc analyses of 2026-10-30 (P1--P9) and 2026-11-03 (P10--P23) have their definitions and their own dated deviations in
\texttt{posthoc/README.md}. Their numbers are reported in the manuscript next to the sealed numbers they qualify and in S4 Table; none replaces a
sealed number or changes a preregistered verdict.

\section{Reference solutions}
\label{sec:refs}
The reference is SciPy's DOP853 at relative and absolute tolerances of $10^{-13}$, maximum step 0.01 TU, with dense output.
\begin{itemize}
\item Study B: a Radau IIA solution at $10^{-12}$ agreed with the reference to at most 0.003 km on 600 usable trajectories, and the event type and
time agreed on all of 300 terminating trajectories.
\item Every trajectory of Studies B--D carries a second DOP853 solution ($10^{-12}$, maximum step 0.005 TU). Post hoc (P18), the largest distance
between the two final positions over all usable trajectories was 2.01 km on the Study B confirmation set (one trajectory above 1 km, none above
10 km; median $1.4\times10^{-8}$ km), 0.050 km on the Study C confirmation set, 0.056 km on the Study C lock set, 0.135, 0.115 and 0.253 km on the
Study D confirmation, replication and lock sets: every one at least 50 times below the 100 km tolerance, and outside Study B at least 390 times.
\item Study D (D12, exploratory): the 163 trajectories with the largest $\lambda T$ ($\ge 10.8$, the top 1\%) were re-integrated with DOP853 at the
tightest tolerance SciPy accepts (a requested $10^{-14}$ is raised to machine precision) with steps of at most 0.002 TU, and with LSODA at $10^{-13}$;
the largest disagreements with the reference were 0.007 km and 0.23 km. No label flipped. The policy violated on 41/163 (25.2\%) of these
trajectories, so its risk in the chaotic tail is its own error, not the reference's.
\item Periodic orbits (D5): a run was excluded when a second reference (DOP853, $10^{-12}$, step at most 0.004 TU) disagreed with the first by more than
5 km (102 of 1,749 runs).
\end{itemize}

\section{Step-size labels and non-monotone labels}
\label{sec:labels}
The label is the largest step on a geometric grid (ratio $2^{1/4}$) that passes together with its eight finer neighbours, refined by up to six
points towards the next coarser grid value while every finer refinement point still passes; the physical-time grid starts at $\min(T, 2\,\mathrm{TU})$,
the fictitious-time grid spans $[10^{-4}, 4]$. A label is censored when the largest grid value passes.
\begin{itemize}
\item Study A (H4): 3 of 200 trajectories (1.5\%; 95\% interval 0.3--4.3\%) had a non-monotone pass/fail sequence on a dense sweep of the physical RK4
step; the median relative disagreement of the halving search with the dense sweep was 0.82\%, and of the monotone label 0.80\%.
\item Study C pilot (Sundman time): 9 of 200 (4.5\%; 2.1--8.4\%; 6\% at $\tau_c = 1$), above the plan's acceptance criterion of $<2\%$ (a recorded
deviation, Section~\ref{sec:dev}); 2.0\% had a failing step inside the label window.
\item Study D (D10, exploratory): 20 of 500 (4.0\%; 2.5--6.1\%). The non-monotone labels were concentrated at short horizons (8\% at 1--3 TU, none at
8--20 TU), far from the Moon (7\% beyond 0.8 DU against 1\% below 0.05 DU) and at low $\lambda T$ (median 2.5 against 3.9), not near collisions.
The suspicion that the landing steps cause them is only weakly supported: the number of physical-time landing steps changed by 0.55 at a
pass/fail flip against 0.49 between any two neighbouring steps.
\end{itemize}
\textbf{Why non-monotone labels do not threaten the certificate.} The label is only the training target. The safety scale is certified on
end-to-end runs of the whole policy (the predicted step times the scale), and each calibration run is judged by its own final error, so a
mislabelled trajectory can make the policy more conservative or more expensive but cannot make the certified violation rate wrong.

\section{The Sundman-time integrator}
\label{sec:sundman}
\textbf{Landing.} Before each fictitious step the predicted physical step $g(\boldsymbol{X}_n)\Delta s$ is compared with the remaining time. If
$t_n + 1.5\,g\Delta s \ge T$, the run finishes with $n_f = \lceil (T-t_n)/(g\Delta s)\rceil$ classical RK4 steps of size $(T-t_n)/n_f$ in physical
time, so the last steps are no longer than the local transformed step and no root finding is needed. A fictitious step whose realized $t$ would
pass $T$ is discarded (its four evaluations stay counted) and the run lands the same way. Every step, fictitious or landing, costs four RHS
evaluations.

\textbf{Order.} On a vertical $L_2$ periodic orbit (stability index 1.0, one period, $T = 6.3$) the Sundman-time RK4 converged at fourth order to
the reference as the step shrank (673 km at $\Delta s = 0.1$, 3.9 km at 0.03, 0.04 km at 0.01); a unit test checks the order on random states.

\textbf{The cap $\tau_c$.} In the Study C pilot the oracle fixed Sundman step was cheaper than the oracle-tuned RK45 on 39.6\%, 45.9\% and 49.5\% of
trajectories for $\tau_c = 0.5$, 1 and 2, and $\tau_c = 2$ was chosen. In the exploratory D10 sweep on 1,500 trajectories of the Study D replication set
($n = 1{,}136$ with an oracle-tuned RK45 cost) the shares were 47.5\%, 51.6\%, 53.2\% and 54.0\% for $\tau_c = 1$, 2, 4 and 8 (95\% intervals about
$\pm3$ points). The narrow threshold result of C3(a) (51.1\%, and 50.6\% on replication) therefore depends on $\tau_c$.

\section{Periapsis features}
\label{sec:peri}
For body $i$ ($\mathrm{GM}_1 = 1-\mu$ at $(-\mu,0,0)$, $\mathrm{GM}_2 = \mu$ at $(1-\mu,0,0)$) let $\boldsymbol{r}$ be the position relative to the body and
$\boldsymbol{v} = (\dot x - r_y,\ \dot y + r_x,\ \dot z)$ the inertial velocity relative to it. With $\boldsymbol{h} = \boldsymbol{r}\times\boldsymbol{v}$,
$\boldsymbol{e} = \boldsymbol{v}\times\boldsymbol{h}/\mathrm{GM} - \boldsymbol{r}/|\boldsymbol{r}|$ and $e = |\boldsymbol{e}|$, the periapsis radius is
$r_p = |\boldsymbol{h}|^2/(\mathrm{GM}(1+e))$, floored at $10^{-8}$ DU. The time to the next periapsis $t_p$ follows from the mean anomaly (elliptic) or the
hyperbolic anomaly (hyperbolic) as in the manuscript, Eq~(peri). Edge cases:
\begin{itemize}
\item near-circular orbits, $e < 10^{-8}$: $t_p = 0$ and $r_p = |\boldsymbol{r}|$ (flag 1);
\item near-parabolic orbits, $|1-e| < 10^{-6}$: the time follows from Barker's equation, and a receding parabola has no future periapsis (flag 2);
\item an outgoing hyperbola has no future periapsis ($t_p = \infty$).
\end{itemize}
Per body the features are $\log_{10} r_p$, $\log_{10}\tau_p$ with $\tau_p = \sqrt{r_p^3/\mathrm{GM}}$, $\min(t_p/T, 2)$ (2 meaning beyond the horizon or
no future periapsis), the indicator of $t_p \le T$, the logarithm of the current distance (relative to the Hill radius $(\mu/3)^{1/3}$ for the Moon) and
$e$; the thirteenth feature is $\log_{10}$ of the smallest over the two bodies of $\tau_p$ if the periapsis is reached before $T$, else
$\sqrt{|\boldsymbol{r}|^3/\mathrm{GM}}$. The flags are not model inputs. The features are exact for two-body motion only. The free crash flag
used by the D4 strata (osculating periapsis inside a collision radius, reached before $T$) had a recall of 81.1\% and a precision of 96.3\% for the
terminating trajectories of the Study C confirmation set and fired on 0.73\% of usable ones.

\section{Certification}
\label{sec:cert}
\textbf{Fixed sequence (Studies C and D).} For the grid $s_1 < \dots < s_{28}$ ($s_k = 0.1\cdot2^{(k-1)/8}$, $k \le 27$, and $s_{28} = 1$), each value is
run end to end on the $n$ calibration trajectories in ascending order; with $V(s_k)$ violations, $p(s_k) = \Pr[\mathrm{Bin}(n, 0.05) \le V(s_k)]$; the
sequence stops at the first $p(s_k) > \delta = 0.1$ and the certified scale is the last value before it. Calibration violations include runs that
collided, escaped, were non-finite, hit the step cap, or abstained to a fallback solver that then violated.

\textbf{Bonferroni (Studies A and B).} Each of 18 values between 0.02 and 1 tested alone at level $\delta/18$; the largest certified value kept. The
adaptive solvers' certified tolerance used the same rule on 18 log-spaced relative tolerances from $10^{-12}$ to $10^{-3}$ in every study.

\textbf{Equal risk.} The loosest relative tolerance on the 18-value grid whose calibration violation rate does not exceed the policy's calibration
violation rate, scanning from the loosest value, refined by four bisections in $\log_{10}$ towards the next looser grid value.

\textbf{Symmetric certification (post hoc, P7).} The solvers certified by the policy's rule (fixed sequence, tightest first, full $\delta$) on
$10^{-10+k/8}$, $k = 0,\dots,56$; see the manuscript, Post hoc checks.

\textbf{Strata of D4.} Base cells: regime group (I; II--III; IV--V) $\times$ lunar class (predicted collision by the free crash flag; osculating lunar
periapsis below 0.05 DU reached before $T$; 0.05--0.17 DU reached before $T$; far) $\times$ horizon band (1--3, 3--8, 8--20 TU), 36 cells. A base
cell with fewer than 400 calibration trajectories first joins the other small cells of the same regime group and lunar class along the horizon
axis; if that union is still below 400 it joins the small remainder of the same regime group along the lunar-class axis; anything still below 400
is pooled. The rule is fixed in advance and applied to the calibration counts. A stratum with fewer than 150 calibration trajectories, or without
a certified scale, would have abstained to RK45 at the certified tolerance (none did). Each stratum certifies its scale by the fixed sequence at
level $\delta/G$ ($G = 17$). Strata with at least 300 evaluation trajectories enter the hypothesis; each is tested by the exact one-sided binomial
test of a rate of at least 7\%, and the intersection--union p-value is the largest. At a 5\% level the same strata give $p = 0.021$ (post hoc, P16).

\section{The leapfrog integrator}
\label{sec:leapfrog}
A finite-difference test of $\Phi^{\mathsf T} J\Phi - J$ in the canonical variables of the rotating frame on 100 states gave a median defect of
$1.1\times10^{-10}$ for the leapfrog step at every step size, the finite-difference floor; RK4's defect was $5\times10^{-5}$ at $\Delta t = 0.2$ and
$4\times10^{-8}$ at 0.05. Leapfrog was dropped in Study C after a comparison on development data: RK4 alone (scale 0.80,
violations 4.5\%, median integration cost 472) against RK4 with a gated leapfrog and one common scale (scale 0.734, violations 3.8\%, median
integration cost 520). Once the certificate covered the whole policy, the leapfrog's cheaper step did not pay for the lower common scale it forced.

\section{Terminating trajectories}
\label{sec:term}
The 3,808 terminating trajectories of the Study C confirmation set (19.0\% of 20,000; the reference collides or escapes before $T$) were run exactly
like usable ones in a separate ledger. The Sundman-time policy ended in the same event type as the reference on 3,805/3,808 (99.9\%: all of 237
escapes, 99.8\% of 600 collisions with the Earth, all of 2,971 collisions with the Moon) and never ran through a body. The event time agreed with the
reference's within $\tol/|\boldsymbol v|$ (the strict rule of the plan) on only 46.7\% (Moon collisions 56.7\%, Earth 15.2\%, escapes 0.4\%), and the
reference event lay inside the detecting step on 96.8\%. The physical-time policy's abstention rule sent 85.4\% of these trajectories to RK45; of the
555 it ran, the event type matched on 93.3\%, but only on 25.5\% of the 47 Moon collisions (51\% were never detected). The certificate of the
manuscript is therefore conditional on survival to $T$.

\section{Endpoint-only audit}
\label{sec:audit}
1,000 Sundman-time runs of the Study C confirmation set were re-run with the error checked at 50 checkpoints on the dense reference. On 2.6\%
(1.7--3.8\%) the largest checkpoint error exceeded the tolerance while the final error did not; the largest error exceeded the tolerance on 7.7\% of
the audited runs against 5.1\% for the final error alone, and nine merged cells exceeded the preregistered 2\% rule. For the physical-time policy the
first share was 0.5\% (0.2--1.2\%). A run judged by its final position can be wrong in mid-run 2--3 points more often than the certificate states.

\section{Exploratory studies of Study B}
\label{sec:studyb}
\textbf{Lookahead (E3c).} An RK45 solve over the first fraction $f$ of the horizon supplied closest approaches and timescales as features (Study B
exploration set, RK4 only, scale certified by the sealed rule). The median error of the predicted step fell steadily with $f$, from 0.32 dex
($f = 0$) to 0.07 dex ($f = 1$), with no early plateau. Lookahead plus integration was cheaper than the certified RK45 on 27.0\% of trajectories at
$f = 0$, 29.5\% at $f = 0.05$ and 6.9\% at $f = 1$, and below 1\% against the oracle-tuned RK45 at every $f$ (model cost excluded); the lookahead itself
cancels the saving.

\textbf{Mass-ratio dial (E7).} The same study at $\mu = 10^{-6}$, $10^{-4}$, $10^{-3}$, 0.01215 and 0.05 (8,000 trajectories each, sampling scaled by the
Hill radius, tolerance 100 km in Earth--Moon units). The median paradigm gap rose with $\mu$ (+0.06, +0.10, +0.11, +0.17 and +0.21 dex) together
with the median timescale range (1.53 to 2.62), and the slope of the gap on $\log_{10}\Delta_\tau$ was 0.71--0.93; the violation rate (2.7--3.7\%) did
not trend. The cross-mass-ratio transfer question was not tested.

\section{Secondary results of Study A}
\label{sec:studya}
\begin{itemize}
\item H2: certified scales gave 17/500 (3.4\%; 1.99--5.39\%) violations for RK4 and 19/498 (3.8\%; 2.31--5.89\%) for leapfrog on the locked audit rows,
against 11.2\% and 14.5\% without the scale (Euler 27.1\%). Upper bounds above 5\% are consistent with a certificate that allows the true rate to
exceed 5\% with probability up to 0.1.
\item H3: removing FTLE changed the RK4 $R^2$ by $-0.25$ ($-0.54$ to 0.04) with the pinned model and by $-0.07$ ($-0.14$ to 0.002) with a random forest,
and selection accuracy by 0.008 ($-0.009$ to 0.024) and $-0.001$ ($-0.011$ to 0.009). Every interval included zero; the point estimates are losses.
\item H4: labels, Section~\ref{sec:labels}.
\item H5: the conformal error-model recommender violated on 15 of 86 covered test trajectories (17.4\%; 10.1--27.1\%), above its 10\% target, and was
never cheaper than RK45 or DOP853.
\item H6: crash or escape prediction with FTLE reached an AUC of 0.993 (recall 86.1\%) and without it 0.780 (recall 19.9\%); the FTLE is integrated over
the same horizon as the event it predicts.
\item H7: median fitted convergence slopes were 1.000 (Euler), 2.000 (leapfrog) and 4.017 (RK4).
\item H1$'$b (post hoc comparator): the model was cheaper than the best single certified fixed RK4 step on 89.2\% and than the best fixed leapfrog step on
98.1\% of trajectories.
\item Sun--Earth ($\mu = 3.00348\times10^{-6}$, 5,000 trajectories, tolerance 38,917 km in the same normalized units): integrator selection reached
91.8\% (89.2--94.4\%) against 61.1\% (52.1--70.0\%) for always choosing RK4, within that system. Transfer between the systems did not beat always-RK4
in either direction (Earth--Moon to Sun--Earth 61.8\% against 61.1\%; Sun--Earth to Earth--Moon 81.9\% against 86.1\%).
\end{itemize}

\section{IAS15 (exploratory)}
\label{sec:ias15}
REBOUND 5.2.2 (inertial frame, the Pham--Rein--Spiegel adaptive step criterion) on 300 evaluation trajectories of the Study D confirmation set, at
the loosest accuracy parameter $\epsilon$ that met 100 km per trajectory, searched up to $\epsilon = 0.1$. The median chosen $\epsilon$ was the edge of the
search ($10^{-1}$), so the comparison is not at a converged setting. A run took a median of 217 $\mu$s (11 steps) against 6.4 $\mu$s for the compiled
policy and 9.2 $\mu$s for compiled DP5(4) on the same trajectories, never faster on a trajectory. The IAS15 times include REBOUND's per-call overhead
from Python, and REBOUND exposes no force-evaluation counter, so the comparison is wall time only and bounds the algorithmic gap from above. IAS15 is
designed for machine-precision integration; we report this as an answer to the question of whether it was tried, not as a competitor, and draw no
conclusion from it in the manuscript.

\section{Compiled timing: the kernel chronology}
\label{sec:d2}
On the Study D selection set, before the seal: (1) the first dry run used the Study C integrator kernel for the policy, which took 26.5 ns per
RHS-equivalent against 18.6 ns for compiled DP5(4), because the Sundman right-hand side evaluated extra square roots and the landing steps allocated
arrays (throughput ratio 0.58); (2) a micro-optimized policy kernel together with a run-time switch in the comparators' right-hand side made the
\emph{comparators} 50\% slower through a dispatch layer, and its ratio of 1.04 was discarded; (3) the reported version compiles both arms with the
right-hand side as a compile-time constant (policy 22.6 ns, DP5(4)-PI 23.1 ns per RHS-equivalent in the dry run). The confirmatory D2 step ran once,
after the seal, with version (3).

\section{Literature search}
\label{sec:lit}
The literature overview (\texttt{LITERATURE\_REVIEW.md}) was compiled on 2026-10-29 from a web search, with each item marked as verified or
unverified against the publisher's page. The databases and queries were not recorded, so the search cannot be reproduced exactly; the statement in
the manuscript that we are not aware of an earlier study of this question rests on that search and on the sources cited there.

\section{Index of the post hoc analyses}
\label{sec:posthoc}
\begin{tabular}{|l|L{5.4in}|}
\hline
P1--P3 & ratio of totals, mean log ratio and share by horizon, regime and family; declared mixtures; one overhead on every column \\ \hline
P4 & three quarters on one dataset \\ \hline
P5 & the empty cell of Fig~4 \\ \hline
P6 & fit-seed variability \\ \hline
P7, P8 & symmetric comparator certification; absolute tolerance and $\delta = 0.05$ \\ \hline
P10 & realized comparator risk on the evaluation trajectories, McNemar tests, violation severity \\ \hline
P11 & same-subset decomposition of the cost gap \\ \hline
P12 & Study B at other overheads; the Sundman policy at multiples of its overhead \\ \hline
P13 & concentration of cost \\ \hline
P14 & timescale range and horizon (partial correlation) \\ \hline
P15 & share and ratio of totals as a function of the near-Moon weight \\ \hline
P16 & Holm over all 15 tests; D4 at 5\%; resampling floors \\ \hline
P17 & BCa intervals for the ratio of totals \\ \hline
P18 & agreement of the two references \\ \hline
P19 & distance of periodic orbits and section grids to the training set \\ \hline
P20--P23 & exploratory comparators: adaptive solvers in Sundman time, an adaptive solver with a learned tolerance, LSODA, and a fixed-step
eighth-order method in Sundman time \\ \hline
\end{tabular}

\end{document}
